\subsection{Small-Lesion Segmentation as a Sparse-Signal Learning Problem}
\label{sec:problem}

Let $\Omega \subset \mathbb{Z}^{3}$ denote the voxel domain of a 3D MRI
volume. For each voxel $i \in \Omega$, let $g_i \in \{0,1\}$ be the
ground-truth label and $p_i \in [0,1]$ the foreground probability predicted
by the segmentation network $f_{\theta}$. The set of lesion voxels is
$G = \{ i \in \Omega \mid g_i = 1 \}$, and it decomposes into $K$ disjoint
connected components, each corresponding to one lesion:
\begin{equation}
G = \bigcup_{k=1}^{K} C_k,
\qquad
C_k \cap C_j = \varnothing \quad \text{for } k \neq j ,
\end{equation}
where $|C_k|$ is the number of voxels in lesion $C_k$. In multiple
sclerosis the lesion signal is sparse in two ways. Lesions occupy a very
small part of the volume, $|G| \ll |\Omega|$, and their sizes are highly
unequal: lesions with an area below 10~mm$^2$ account for nearly 20\% of
all lesions but for about 1\% of the total lesion load on
average~\cite{wang_survey_1997}.

A voxel-wise loss such as cross-entropy is a sum of per-voxel terms
$\ell(p_i, g_i)$, which can be grouped by lesion:
\begin{equation}
\mathcal{L}_{\mathrm{vox}}
=
\sum_{i \in \Omega} \ell(p_i, g_i)
=
\sum_{k=1}^{K} \sum_{i \in C_k} \ell(p_i, 1)
+
\sum_{i \notin G} \ell(p_i, 0) .
\label{eq:sl_voxel_loss}
\end{equation}
Lesion $C_k$ enters the loss, and therefore the gradient, through $|C_k|$
terms, so its influence on training is proportional to its size.
Background and large lesions dominate, and a lesion of a few voxels barely
changes the objective~\cite{kofler_blob_2023}. Overlap-based losses behave
in the same way and are additionally unstable on small
objects~\cite{azad_loss_2023}. For example, if a prediction is perfect
except that one lesion $C_m$ is missed completely, the DSC is
\begin{equation}
\mathrm{DSC}
=
\frac{2\,(|G| - |C_m|)}{2|G| - |C_m|}
=
1 - \frac{|C_m|}{2|G| - |C_m|} .
\label{eq:sl_missed_lesion}
\end{equation}
A lesion that carries 1\% of the lesion volume lowers the DSC only to
about 0.995, although one of the $K$ lesions has not been detected at all.
Voxel-level overlap therefore does not control lesion-level detection.

The question addressed in this work is: \textit{how can the training
objective be redesigned so that small lesions contribute to optimization as
much as large ones, and so that every lesion is detected?} These are two
separate requirements. The first concerns \emph{balance}: the weight of
lesion $C_k$ in the loss should not scale with $|C_k|$. The second concerns
\emph{detection}: the complete omission of a lesion should be penalized
directly, at the level of the lesion rather than of its voxels. Both are
imposed on the loss only; the network $f_{\theta}$ is left unchanged, so
that differences in performance can be attributed to the objective.

\subsection{CATMIL Objective}

The proposed objective, CATMIL, adds two terms to the standard nnU-Net
loss, one for each requirement: a Component-Adaptive Tversky (CAT) term to
balance lesions and a lesion-level MIL term for detection. The connected
components $C_1,\ldots,C_K$ of the ground-truth foreground are computed
from the ground-truth mask of each training patch using 6-connectivity,
consistent with the definition of a lesion used for evaluation. They are
therefore fixed with respect to the prediction.

\subsubsection{Structure-Level Component-Adaptive Tversky Term}

According to \cref{eq:sl_voxel_loss}, a voxel-wise loss gives lesion $C_k$
an influence proportional to its size $|C_k|$. The CAT term removes this
dependence by weighting each voxel according to the component it belongs
to:
\begin{equation}
w_i =
\begin{cases}
(|C_k| + \varepsilon_{c})^{-\gamma}, & \text{if } i \in C_k, \\
w_{\mathrm{bg}}, & \text{if } g_i = 0 .
\end{cases}
\label{eq:sl_cat_weights}
\end{equation}
The exponent $\gamma \geq 0$ controls the degree of size compensation.
When $\gamma=0$, the weighting remains proportional to the number of
voxels. With $\gamma=1$, as used in this work, the total weight of lesion
$C_k$ is
\begin{equation}
\rho_k
=
\sum_{i \in C_k} w_i
=
\frac{|C_k|}{|C_k| + \varepsilon_{c}} ,
\label{eq:sl_cat_lesion_weight}
\end{equation}
which is bounded by one and no longer increases linearly with lesion
volume. The offset $\varepsilon_{c}=5$ voxels limits the weight assigned to
each voxel, preventing lesions consisting of only a few voxels from
receiving excessively large per-voxel weights. Thus, $\rho_k$ is $0.67$ for
a 10-voxel lesion and $0.95$ for a 100-voxel lesion. Since
$\rho_k \in [1/6, 1)$ for any lesion of at least one voxel, lesion volumes
may differ by orders of magnitude while their total weights differ by less
than a factor of six. Because the foreground weights are much less than
one, the background weight is set to $w_{\mathrm{bg}}=0.01$ so that the
false-positive term remains on a comparable scale.

With these weights, the true-positive, false-positive, and false-negative
sums are
\begin{equation}
\mathrm{TP}_{w} = \sum_{i \in \Omega} w_i\, p_i\, g_i,
\qquad
\mathrm{FP}_{w} = \sum_{i \in \Omega} w_i\, p_i\, (1-g_i),
\qquad
\mathrm{FN}_{w} = \sum_{i \in \Omega} w_i\, (1-p_i)\, g_i ,
\end{equation}
and the CAT term is one minus the corresponding Tversky
index~\cite{salehi_tversky_2017}:
\begin{equation}
\mathcal{L}_{\mathrm{CAT}}
=
1 -
\frac{\mathrm{TP}_{w}+\delta}
{\mathrm{TP}_{w} + \alpha\, \mathrm{FP}_{w} + \beta\, \mathrm{FN}_{w} + \delta},
\label{eq:sl_cat_loss}
\end{equation}
where $\alpha=0.3$ and $\beta=0.7$ weight false positives and false
negatives, respectively, as in the Tversky baseline. The parameter
$\delta=10^{-5}$ is the usual smoothing constant. The term is set to zero
when the training batch contains no lesions.

Under these weights, the sums take a simple form. Let
$r_k = \frac{1}{|C_k|}\sum_{i \in C_k} p_i \in [0,1]$ be the mean predicted
probability inside lesion $C_k$, that is, its soft recall. Then
\begin{equation}
\mathrm{TP}_{w} = \sum_{k=1}^{K} \rho_k\, r_k,
\qquad
\mathrm{FN}_{w} = \sum_{k=1}^{K} \rho_k\, (1-r_k) .
\label{eq:sl_cat_lesion_form}
\end{equation}
Lesions therefore contribute to the loss through their recalls $r_k$, with
coefficients $\rho_k$ that are nearly independent of lesion size. A
completely missed lesion ($r_m=0$) transfers a total weight of
$\rho_m < 1$ from $\mathrm{TP}_{w}$ to $\mathrm{FN}_{w}$, whereas in
\cref{eq:sl_missed_lesion} its effect on the DSC is proportional to
$|C_m|$. The gradient is balanced in the same way: for a voxel
$j \in C_k$,
$\partial \mathcal{L}_{\mathrm{CAT}} / \partial p_j = -c\, w_j$ with a factor
$c>0$ common to all foreground voxels (\cref{appendix:catmil-properties}),
so the total gradient magnitude received by lesion $C_k$ is $c\,\rho_k$
instead of increasing with $|C_k|$. For $\gamma=0$ and $w_{\mathrm{bg}}=1$,
the term reduces to the standard Tversky loss; if additionally
$\alpha=\beta=0.5$, it is the soft Dice loss. The CAT term can therefore be
viewed as a Tversky loss in which the lesion, rather than the voxel, is the
effective unit of counting.

\subsubsection{Lesion-Level Multiple Instance Learning Term}

The CAT term balances lesions but still evaluates them by overlap: it is
linear in $r_k$ and does not distinguish between a partially segmented
lesion and one that is not detected at all. The second requirement is
addressed using the Multiple Instance Learning
paradigm~\cite{ilse_attention-based_2018}. Each lesion is treated as a bag
of voxels. A bag is considered detected if at least one of its voxels
receives a high foreground probability. The detection score of lesion
$C_k$ is therefore its strongest response,
\begin{equation}
s_k = \max_{i \in C_k} p_i ,
\end{equation}
and the MIL term is the average negative log-score over lesions:
\begin{equation}
\mathcal{L}_{\mathrm{MIL}}
=
-\frac{1}{K} \sum_{k=1}^{K} \log \max(s_k, \varepsilon) ,
\label{eq:sl_mil_loss}
\end{equation}
where $\varepsilon=10^{-6}$ keeps each contribution finite and $K$ denotes
the number of lesions in the training batch. If the batch contains no
lesions, the term is set to zero.

This term directly penalizes the complete omission of a lesion. If lesion
$C_k$ is missed, all $p_i$ in $C_k$ are close to zero, so $s_k \approx 0$
and its contribution, $-\log s_k$, is large. Once a single voxel inside the
lesion has a high predicted probability, $s_k$ approaches one and the
contribution vanishes. The gradient acts on the maximal voxel of each
lesion with magnitude $1/(K s_k)$ (\cref{appendix:catmil-properties}).
It does not depend on $|C_k|$ and is largest for lesions that are currently
missed. Consequently, a lesion consisting of only a few voxels receives the
same detection signal as a large lesion. The term does not require the
entire lesion to be segmented; delineation is left to the overlap terms.

\subsubsection{Objective Function}

The complete objective adds the two terms to the default nnU-Net loss. In
this binary setting, the default loss is the sum of the Dice and
cross-entropy losses,
$\mathcal{L}_{\mathrm{base}} = \mathcal{L}_{\mathrm{Dice}} + \mathcal{L}_{\mathrm{CE}}$:
\begin{equation}
\mathcal{L}_{\mathrm{CATMIL}}
=
\mathcal{L}_{\mathrm{base}}
+
\lambda_{\mathrm{CAT}}(t)\, \mathcal{L}_{\mathrm{CAT}}
+
\lambda_{\mathrm{MIL}}\, \mathcal{L}_{\mathrm{MIL}} .
\label{eq:sl_catmil_loss}
\end{equation}
The base loss provides voxel-level classification and global overlap, the
CAT term provides component-balanced overlap, and the MIL term provides
lesion-level detection. Only the loss is modified; the network $f_{\theta}$
and, therefore, its hypothesis class remain unchanged. Both auxiliary terms
are bounded and continuous in the predictions
(\cref{appendix:catmil-properties}), so their influence is controlled by
the two coefficients. The base loss retains nnU-Net's deep supervision,
whereas the CAT and MIL terms are computed only on the full-resolution
output.

\paragraph{Training strategy.}
The MIL term is assigned a constant weight of
$\lambda_{\mathrm{MIL}}=0.2$. The CAT term is introduced using a linear
warm-up:
\begin{equation}
\lambda_{\mathrm{CAT}}(t)
=
\lambda_{\mathrm{CAT}} \cdot \min\left(\frac{t}{T},\, 1\right),
\qquad
\lambda_{\mathrm{CAT}} = 0.1 ,
\end{equation}
where $t$ is the current training epoch and $T=50$ epochs is the warm-up
duration. Because the CAT term changes how lesions contribute to the loss,
applying it at full strength from the first epoch may destabilize
optimization before the model has learned basic foreground-background
separation. The warm-up allows the model to first establish stable
voxel-level segmentation and then gradually shifts the supervision toward
component balance. The weights $\lambda_{\mathrm{CAT}}$ and
$\lambda_{\mathrm{MIL}}$ were fixed before the experiments and are used
unchanged for both datasets; their effect is examined in a sensitivity
analysis in \cref{appendix:catmil-loss-weights}.
